\begin{table}[htbp]
\centering
\begin{threeparttable}
\caption{Production Caps and Seismic Activity: The Mechanism}
\label{tab:mechanism}
\begin{tabular}{lcccc}
\toprule
 & \multicolumn{2}{c}{Production (bcm)} & \multicolumn{2}{c}{Earthquakes ($M \geq 1.0$)} \\
\cmidrule(lr){2-3} \cmidrule(lr){4-5}
Period & Mean & Range & Annual Mean & Max Magnitude \\
\midrule
Pre-cap (2003--2013) & 45.7 & 36.2--53.9 & 42.1 & 3.6 \\
Cap era (2014--2019) & 27.0 & 18.9--42.5 & 43.7 & 3.4 \\
Wind-down (2020--2023) & 6.2 & 0.0--11.8 & 28.0 & 3.2 \\
\bottomrule
\end{tabular}
\begin{tablenotes}[flushleft]
\small
\item \textit{Notes:} Production data from NAM/Rijksoverheid official publications. Earthquake counts from KNMI FDSN catalog (magnitude $\geq 1.0$, Groningen region: lat 52.5--53.8, lon 5.5--7.5). The production-to-earthquake link is physically determined: gas extraction causes compaction-driven fault slip. The decline in both production and seismicity after 2014 reflects the government's production cap policy.
\end{tablenotes}
\end{threeparttable}
\end{table}
